% ==========================================================================================
\section{Discussion}
\label{sec:discussion}

\paragraph{Disclosure reaches the operators who were already upgrading}
The LND advisories of 2025 and 2026 were published long after the fixes shipped, a common practice that lets
operators upgrade before attackers learn what to look for. It works for the operators who follow releases: half of
LND nodes move to a new line within three months, and the largest move first. The nodes still on an affected line at
disclosure are the ones that did not, and the disclosure itself did not change their behaviour. Their upgrade rate in
the month after Lightning Labs' batch was the same as in the month before, and most of them were still exposed two
months later. Once a release is a year old, the exposed tail consists of nodes that upgrade rarely or never, which an
advisory alone does not reach. CLN's July disclosures did move operators: the rate at which nodes left vulnerable
lines rose two- to fourfold in the following month, before the maintainers' call of 26~August, after which exposure
fell further. We cannot separate the reasons. The fixed CLN lines were younger (one of the two fixes had shipped \FixToDiscClnGossipd{}
days before disclosure, against \FixToDiscLLBatch{} for LND), so more of the exposed operators were ones who usually keep up; CLN's
operators are fewer and closer to its maintainers; and its disclosures came amid an active run of CLN security
releases. Part of CLN's decline was nodes going offline rather than upgrading, which the maintainers asked for and
which also removes the exposure.

\paragraph{Platforms move nodes; size and attention keep them current}
Upgrades arrive in waves of a hundred or more nodes a few weeks after a release, consistent with node platforms
shipping updates to all their users at once, and one wave followed a merchant platform's security release within a
day. This matches what Frei et al.\ found for browsers: the update mechanism, more than the user, sets the patch
level~\cite{frei2009firefox}. The tail is the opposite population: small nodes, nodes that skipped earlier releases,
nodes no platform updates for them. For advisories that require a channel with the victim, these nodes are also the
cheapest to reach. Their share of capacity at risk is small (about \ExpCapLLBatchDisc\% for the August~2026 LND
batch), but each holds its operator's funds.

\paragraph{Counting nodes}
For half a year, one operator ran more than a fifth of the live LND nodes we observe, with almost no capacity. Any
statistic that counts nodes, whether implementation shares, version shares or exposure, would have been dominated by
that operator's choices. Fleets are cheap to create and hard to see; we found this one only because it upgraded in
step. Node-level Lightning measurements should look for such fleets before counting, and report capacity alongside
nodes~\cite{kumar2026clientdiversity}.

% ------------------------------------------------------------------------------------------
\section{Limitations}
\label{sec:limitations}
\emph{Release lines, not patches.} Gossip cannot separate patch releases, so fixes inside a line, including all
of CLN's 2026 patch releases, cannot be measured; for advisories with a fix inside a line we report bounds.
\emph{Inference.} An operator can switch features on or off so that an older line looks like a newer one or the
reverse; the strict range covers every non-developer configuration, and validation on self-identifying CLN nodes
bounds the error of the likely line, but we have no equivalent ground truth at scale for LND. A node that upgrades
without re-announcing, or whose announcement we do not receive, keeps its old bits, which biases exposure up.
\emph{Vantage and liveness.} All data come from one node's view of the public graph; unannounced nodes, including
most mobile wallets, are not covered, and liveness is inferred from peers' channel updates.
\emph{Fleets.} We identified one managed fleet by its alias and hubs. Smaller fleets without a shared alias would go
undetected and be counted as independent nodes.
\emph{Exposure is not exploitability.} Several advisories require a channel with the victim, a particular
configuration or inbound connections; we count nodes on affected releases, not nodes that can be attacked. Our
advisory list is compiled from public sources and may miss silent fixes.
\emph{Association.} The hazard model describes which nodes upgrade sooner; it does not identify why.

% ------------------------------------------------------------------------------------------
\section{Ethics}
\label{sec:ethics}
Release-line inference identifies nodes running vulnerable software, and in 2026 Lightning deployments were
exploited (BTCPay Server~\cite{btcpay2026advisory}) and CLN maintainers reported attacks on unpatched nodes. We therefore publish aggregates only: no list, ranking, filter or label of
individual nodes by release or exposure, in the paper or in the released data. All inputs are public gossip,
every advisory we study had been published before our analysis began, and the paper reports no new vulnerability.

% ------------------------------------------------------------------------------------------
\section{Conclusion}
\label{sec:conclusion}
Lightning gossip carries no version string, but the feature bits each release announces identify its release line
closely enough to measure how the network patches. Most operators upgrade within a few months, in waves pushed by
their platforms; a tail of small and neglected nodes does not, and LND's disclosures, which typically come long after
the fix, leave that tail where it was. CLN's 2026 disclosures, followed by the maintainers' appeal, moved theirs. Measured by capacity, the
exposure is small; measured by nodes, it depends on whether one counts a single operator's fleet as hundreds of
operators.
